\section{Extensions and corrections to the framework}
\label{sec:numerical-framework:extensions}

The framework was extended to support the chosen convective flux,
reference-consistent molecular transport and additional convergence
diagnostics. A separate correction addresses the variable-density SA
formulation. The following descriptions identify the purpose and scope of
each change. Their source is supplied through the
\hyperref[chap:data-avail]{Data Availability statement}.

\subsection{HiSA extension with \texorpdfstring{AUSMDV2\(^+\)}{AUSMDV2+}}
\label{subsec:ext:AUSM}
% Source provenance (WSL):
% /home/mgnadlinger/OpenFOAM/mgnadlinger-v2412/modules/hysa/src/finiteVolume/
% fluxSchemes/ausmDV/ausmDVFluxScheme.{H,C}
% TypeName("AUSMDV"); AUSMDV2+ blending with locally modified pressure flux.

The in-house HiSA extension provides an AUSMDV2\(^+\)-based flux scheme,
selected as \texttt{AUSMDV} in the case dictionaries. It blends AUSMD\(^+\)
and AUSMV\(^+\) contributions following
Phongthanapanich and Matthujak~\cite{PhongthanapanichMatthujak2021AUSMDV2Plus}.
The scheme was selected for shock resolution and robustness. The
stand-off-distance study in
Section~\ref{sec:numerical-vv:shock-stand-off-carbuncle} assesses this
implemented variant and documents its remaining carbuncle sensitivity.
The later boundary layer assessment separately examines the viscous
transport operations introduced in
Section~\ref{sec:numerical-framework:discretization}.

\subsection{Additional thermophysical-model combinations}
\label{subsec:ext:thermo-physical-extension}
% Source provenance (WSL):
% /home/mgnadlinger/OpenFOAM/mgnadlinger-v2412/src/myThermos/myThermos.C

Matching the DNS references requires independent control of the molecular
viscosity and conductivity laws. The user library \texttt{myThermos}
registers OpenFOAM's existing polynomial transport model with the ideal-gas
equation of state, sensible internal energy, and either constant or
temperature-dependent specific heat. It enables these combinations without
changing the underlying property models.
\par\medskip

For the calorically perfect benchmarks, seventh-order polynomials represent
the prescribed viscosity laws, and proportional conductivity coefficients
enforce the selected constant molecular Prandtl number. The fits are checked
over the temperatures encountered in each simulation.
Section~\ref{sec:evaluation:benchmark-cases} documents the reference laws
and fit assessment.

\subsection{Density-gradient correction to Spalart--Allmaras}
\label{subsec:ext:spalart-allmaras-of-correction}

The OpenFOAM v2412 SA implementation places density inside the diffusive
flux of its transported working variable. For variable density, this
introduces an additional density-gradient term relative to the
material-form SA equation~\cite{SpalartAllmaras1994,AJS2012}.
The custom \texttt{SpalartAllmarasDensityCorrected} model cancels that term
while retaining the other components of the local implementation.
Appendix~\ref{app:spalart-allmaras-density} gives the derivation and source
audit. This correction restores the intended continuum diffusion structure and is subsequently used as the only SA model being assessed in Section~\ref{sec:comparative:zpg:wall-resolved-baselines}.